% ======================================================================
% ZAKLJUČAK (nenumerisano)
% ======================================================================
\phantomsection
\addcontentsline{toc}{chapter}{Zaključak}
\chapter*{Zaključak}
\label{ch:zakljucak}

Eksperiment proveden kroz 1.070 validnih mjernih serija nad skupom podataka StackOverflow2010 s $\sim$19,2 miliona zapisa pokazuje da HTAP dual-store konfiguracija dominira u svim ispitanim scenarijima mješovitog opterećenja. Pri razini konkurentnosti $C=32$, HTAP postiže 21.920 ops/s OLTP propusnosti, što je 4,6 puta više od IM konfiguracije, 8,8 puta više od CS konfiguracije i 9,0 puta više od Baseline konfiguracije. Najniže OLAP latencije ($P_{50}=924$\,ms pri 80/20 opterećenju), uz relativno malu standardnu devijaciju, svjedoče o stabilnom i predvidivom radu sistema. Ovu dominaciju treba čitati uz ograničenja iz odjeljka~\ref{sec:ogranicenja}: ispitana dual-store konfiguracija odvaja slojeve bez kontinuirane sinhronizacije, pa izmjerene prednosti predstavljaju gornju granicu pristupa.

IM konfiguracija pokazuje iznenađujuću slabost pri mješovitom opterećenju unatoč korištenju istog Hekaton engina kao HTAP: row-mode izvršavanje OLAP upita zasićuje procesorske resurse -- mehanizam izveden iz plana izvršavanja, a ne iz direktnog mjerenja -- i smanjuje OLTP propusnost za faktor 3,3 do 6,9 u odnosu na HTAP. CS konfiguracija bilježi degradaciju OLTP propusnosti od 67,4\% pri prelasku s transakcijskog na mješovito opterećenje, što u apsolutnim vrijednostima znači gubitak dvije trećine OLTP kapaciteta pri $C=32$. Baseline konfiguracija potvrđuje svoju ulogu referentne točke i nije preporučena za produkcijska okruženja s analitičkim zahtjevima.

Izvorni doprinos ovog rada ogleda se u tri aspekta. Kao prvo, identificiran je i kvantificiran fenomen koji dosad nije bio eksplicitno mjeren u literaturi: columnstore ubrzanje OLAP upita indirektno poboljšava OLTP propusnost pri mješovitom opterećenju za faktor 3,3--6,9, jer oslobađa CPU kapacitet koji bi u alternativnoj konfiguraciji (IM) bio potrošen na row-mode skeniranja $\sim$3,7 miliona redova. Ovaj nalaz mijenja način na koji treba pristupiti optimizaciji HTAP sistema: poboljšanje OLAP performansi nije samo cilj samo po sebi, već i sredstvo za poboljšanje OLTP performansi.

Kao drugo, rad pruža prvu sistematičnu komparativnu evaluaciju četiri arhitekturalne konfiguracije SQL Servera 2022 nad stvarnim skupom podataka s Pareto distribucijom korisničke aktivnosti, u punom faktorijalnom eksperimentalnom dizajnu $4 \times 3 \times 3$ s 30 ponavljanja po kombinaciji i rigoroznom neparametarskom statističkom analizom.

Kao treće, empirijski nalazi prevode se u konkretne, akcijske preporuke za inženjere i arhitekte koji projektuju HTAP sisteme u SQL Server okruženju: za sve scenarije s mješovitim opterećenjem, bez obzira na omjer OLTP/OLAP aktivnosti, dual-store konfiguracija dosljedno daje bolje rezultate od alternativa -- uz preduslove navedene u odjeljku~\ref{sec:ogranicenja} (dovoljan RAM i tolerancija na asinhronu svježinu analitičkih podataka).

Rezultati ovog rada otvaraju nekoliko smjernica za buduća istraživanja. Kao prvo, evaluacija na distribuiranim HTAP sistemima (TiDB, SingleStore) s identičnim skupom podataka i workloadom omogućila bi direktnu komparaciju SQL Server in-process dual-store-a s distribuiranim dual-engine arhitekturama i kvantifikaciju troška mrežne latencije. Kao drugo, istraživanje uticaja tipa pohrane (HDD naspram NVMe SSD) na relativne prednosti konfiguracija razjasnilo bi u kojoj mjeri izmjereni trendovi ovise o I/O okruženju, s obzirom da je eksperiment proveden na NVMe SSD-u. Kao treće, evaluacija na novijim Stack Overflow dumpovima s većim skupovima podataka (2015, 2023) pružila bi uvid u skalabilnost izmjerenih trendova. Kao četvrto, istraživanje automatskog tuninga koji dinamički bira optimalnu konfiguraciju ovisno o trenutnom omjeru opterećenja -- ili čak kombinira sve četiri konfiguracije adaptivno -- otvorilo bi put prema samooptimiziranom HTAP sistemu.
